\chapter{Att kontrollera ett påstående}
\label{app:verifiera}

Det här materialet gör, som allt undervisningsmaterial, en rad
påståenden.  De flesta går att pröva på plats: kör programmet och se
efter.  Men några gör det inte.  Att stilreglerna vi följer är
Pythongemenskapens egna, att de verkligen gör koden lättare att läsa för
någon annan, och att de missuppfattningar undervisningen är byggd för att
motverka verkligen är dokumenterade hos nybörjare --- det går inte att se
i en terminal.  Sådana påståenden måste kontrolleras mot källor.

Samma fråga möter er varje gång ni läser en lärobok, en webbsida eller ett
svar från en språkmodell, och varje gång ni själva ska skriva en rapport.
Den här bilagan beskriver en arbetsgång för att kontrollera ett påstående,
och \cref{tab:paastaenden} redovisar var vart och ett av modulens
påståenden är kontrollerat.  \Cref{app:pep8} visar arbetsgången genomförd
på det påstående som krävde en egen sökning: att stilreglerna gör koden
lättare att läsa för någon annan.

\section{Gör påståendet till en fråga}
\label{sec:fraga}

Ett påstående går att kontrollera först när det är tydligt vad som skulle
kunna visa att det är \emph{fel}.  Formulera det därför som en fråga med
ett svar: \enquote{skriver PEP~8 att konstanter sätts med versaler?} kan
besvaras med ja eller nej, medan \enquote{PEP~8 är viktig} inte kan det.
Ingen tänkbar källa skulle visa att det senare är fel, eftersom
\enquote{viktig} är ett omdöme och inte ett faktum.  Den som säger så
menar oftast något annat --- till exempel \enquote{gör en gemensam
stilguide koden lättare att läsa för andra?} --- och det är en fråga som
kan få svaret nej, och därför kan prövas.  Den frågan prövas i
\cref{app:pep8}.

Ofta döljer sig flera påståenden i ett.  \enquote{En dokumenterad
missuppfattning hos nybörjare} är två frågor: finns den beskriven i
litteraturen, och gäller beskrivningen den grupp vi undervisar?  De
kräver olika slags källor.

\section{Välj rätt sorts källa}

Olika frågor besvaras av olika källor.
\begin{description}
  \item[Primärkällan] för vad ett dokument \emph{säger}.  Vill man veta
    vad PEP~8 föreskriver läser man PEP~8, inte en blogg som återger den.
    Att en text finns på webben gör den inte föränderlig i efterhand:
    därför anges hämtningsdatum.
  \item[Enskilda studier] för ett specifikt resultat --- att en viss
    missuppfattning har observerats, hos vilka och hur många.  En enskild
    studie visar att något observerats en gång.
  \item[Översikter] (\foreignlanguage{english}{reviews}) för om något är
    \emph{utbrett} eller vad forskningen \emph{sammantaget} säger.  En
    systematisk översikt går igenom hundratals studier efter en angiven
    metod och väger därför tyngre än en enskild studie.
  \item[Uppslagsverk] för vad ett ord \emph{betyder}.
\end{description}
Var man hittar dem: universitetsbibliotekets söktjänst, och databaser som
Scopus, Web of Science, IEEE Xplore, ACM Digital Library, DBLP (datalogi)
och OpenAlex (öppen; \enquote{öppen} står här för öppna data, inte för
\foreignlanguage{english}{open access}).  Google Scholar täcker brett men
rankar ogenomskinligt.

\section{Sök så att ni kan ha fel}

Den som söker på namn hon redan känner till hittar det hon väntar sig.
Sök därför \emph{ämnesdrivet} --- på begrepp, inte författare --- och ställ
samma frågor i flera databaser.  Sök också efter \emph{invändningar}: lägg
till ord som \enquote{critique}, \enquote{limitations} eller
\enquote{failure}.  Ett påstående som överlevt en motsökning är värt mer
än ett som aldrig prövats, och en invändning som hittas blir ett
\emph{förbehåll} i svaret, inte ett nederlag.

Somliga källor går bara att nå som \emph{kända dokument}: man vet att de
finns och slår upp dem direkt.  Så är det med PEP~8 och PEP~257, som inte
är forskningsartiklar utan gemenskapens egna dokument, och de hämtas
därför från Pythons egen PEP-webbplats.  Det är i sin ordning, bara det
sägs.

\section{Läs i fulltext, och skriv ner hur}

En träff i en databas är inte ett belägg.  Sammanfattningen
(\foreignlanguage{english}{abstract}) räcker inte heller: den säger vad
författarna vill ha sagt, inte vad de visat.  Läs hela texten, leta upp
\emph{det ställe} som styrker påståendet och skriv av det ordagrant, med
sidnummer.  Hittar ni inget sådant ställe stöder källan inte påståendet,
hur passande titeln än är.

Kravet gäller också --- kanske särskilt --- källor ni har ärvt.  En
hänvisning som redan står i ett dokument är inte kontrollerad förrän
\emph{ni} har läst stället den pekar på.  Källorna till modulens
påståenden om missuppfattningar är hämtade ur en artikel om nybörjares
missuppfattningar i programmering\autocite{SooriBosk2026}, och varje
citat som undervisningen vilar på har lästs om i originalkällan innan det
återanvänts här.  Det var inte förgäves: omläsningen visade att tre av
sidhänvisningarna var fel --- citatet stod på en annan sida än
anteckningen påstod, och i ett fall handlade den angivna sidan om något
helt annat.  De är rättade.  Två källor som tidigare bara var lästa i
sammanfattning gick dessutom att få tag på i fulltext, och kunde
kontrolleras ordentligt.

Skriv sedan ner hur ni gick till väga, så att någon annan kan följa
spåret.  I det här materialet står den anteckningen som en kommentar
ovanför varje post i litteraturlistans källfil, med fasta fält:
\begin{description}
  \item[\texttt{CLAIM}] vilket påstående källan ska styrka;
  \item[\texttt{FOUND-VIA}] hur den hittades --- sökfråga och databas,
    eller \enquote{känt dokument};
  \item[\texttt{PICKED}] varför just den, framför andra träffar;
  \item[\texttt{QUOTE}] det ordagranna citatet, med sida;
  \item[\texttt{VERIFIED}] att fulltexten lästs, och var;
  \item[\texttt{COUNTER}] vad motsökningen gav;
  \item[\texttt{DATE}] när.
\end{description}
Anteckningen ser byråkratisk ut och tar ett par minuter per källa.  Den
är också det enda som gör det möjligt att om ett år svara på frågan
\enquote{hur vet du det?}

\section{Dra slutsatsen --- med förbehåll}

Svaret skrivs ut tillsammans med det som talar emot.  Två förbehåll gäller
alla påståenden om missuppfattningar i \cref{tab:paastaenden}.  Det
första: en dokumenterad missuppfattning är dokumenterad hos \emph{några}
studenter, i \emph{några} studier --- inte hos alla, och inte
nödvändigtvis här.  Materialet använder dem därför som skäl att
\emph{visa} en kontrast, inte som påstående om vad just ni tror.  Det
andra: flera av studierna är gjorda i andra språk än Python och i andra
skolformer, vilket gör dem till anledningar att undersöka saken, inte till
bevis.  Frågorna i föreläsningen är därför formulerade så att de kan
besvaras fel utan förlägenhet; svaren säger vad som behöver övas.
Varje kapitel slutar därför med ett avsnitt \emph{Fortsatt arbete}:
vad just den sökningen lämnade ogjort, och vad som skulle behöva
undersökas för ett säkrare svar.

\begin{table}[htbp]
  \centering
  \caption{Modulens påståenden, och var vart och ett är kontrollerat.}
  \label{tab:paastaenden}
  \footnotesize
  \begin{tabular}{@{}p{0.42\linewidth}p{0.5\linewidth}@{}}
    \toprule
    Påstående & Kontrollerat \\
    \midrule
    En tilldelning uppfattas av nybörjare som att variabeln håller hela
      uträkningen, eller som en ekvation datorn ska lösa
      & Återanvänt \parencite{SooriBosk2026}; citaten omlästa i Kohn
        respektive Plass-Oude Bos \\
    En variabel uppfattas kunna hålla flera värden samtidigt
      & Återanvänt; citatet omläst i Sorvas avhandling, tabell~A.1 \\
    Nybörjare har svårt med satsernas följdordning och läser programmet
      som om alla rader vore verksamma samtidigt
      & Återanvänt; citaten omlästa i Sorvas avhandling och i Peas artikel \\
    Inläst numerisk data antas vara tal fast den är en sträng
      & Återanvänt; citatet omläst i Guo, Markel och Zhang \\
    Variabelns namn uppfattas ha makt över dess värde
      & Återanvänt; citatet omläst i Sleeman m.fl., förtecknat i Qian och
        Lehmans översikt \\
    PEP~8 är Pythongemenskapens stilguide och föreskriver versaler för
      konstanter, fyra blanksteg i indrag och högst 79 tecken per rad
      & Primärkälla, hämtad och läst \\
    PEP~257 är konventionen för docstrings
      & Primärkälla, hämtad och läst \\
    Att följa en stilguide som PEP~8 gör koden lättare att läsa för
      någon annan
      & Egen tvåsidig sökning, \cref{app:pep8}: stöd för enskilda regler,
        inte för guiden som helhet \\
    Kontrast måste upplevas samtidigt, och komma före generalisering
      & Primärkälla, Martons \emph{Necessary Conditions of Learning} \\
    \bottomrule
  \end{tabular}
\end{table}

\section{En anmärkning om språkmodeller}

Ett av modulens påståenden har krävt en egen litteratursökning, den i
\cref{app:pep8}; de övriga källorna är antingen primärkällor som slagits
upp direkt eller källor som återanvänts från en forskningsartikel och
lästs om.  I den sökningen --- liksom i de sökningar som gjorts i den här
kursens övriga material, se föreläsningarna \emph{Algoritmiskt tänkande}
och \emph{Funktioner} --- har en språkmodell använts för att
\emph{sortera} träffarna, aldrig för att avgöra om en källa stöder ett
påstående.  Lägg märke till
arbetsfördelningen: modellen sorterar, men varje källa som citeras ska
läsas och väljas av en människa.  Använd gärna språkmodeller för att hitta
och sortera; låt dem aldrig vara det sista ledet.  De kan göra fel, men
det är ni själva som måste stå till svars för innehållet.


\chapter{Hjälper det läsaren att koden följer PEP~8?}
\label{app:pep8}
\chapterprecis{Författaren har ännu inte granskat resultaten i den här
  bilagan i sin helhet.}

Innan den här sökningen stod det i \cref{sec:kommentarer} att ni ska
följa stilreglerna i PEP~8\autocite{pep8} \enquote{av samma skäl som vi
valde beskrivande namn: koden ska gå att läsa för någon annan än den som
skrev den}; det här kapitlet är skälet till att det inte längre står så.
Påståendet som prövas är alltså att koden blir lättare att läsa för någon
annan av att den följer stilreglerna i PEP~8.  Det är inget påstående om
vad PEP~8 \emph{säger} --- det står i PEP~8 och är lätt att slå upp ---
utan ett påstående om vad stilregler \emph{gör} med en läsare.
\Cref{sec:fraga} använder just den frågan som exempel på ett påstående
som går att pröva, och frågan lyder: \emph{blir koden lättare att läsa
för någon annan av att den följer en stilguide som PEP~8, och kostar det
något att låta bli?}

Frågan gäller ett knippe regler, och de behöver skiljas åt.  Fyra
delfrågor prövas var för sig: (a)~gör \emph{namnen} skillnad, (b)~gör
\emph{utseendet} --- indrag, mellanslag, radlängd --- skillnad, (c)~gör
det skillnad att följa en konvention \emph{över huvud taget}, och (d)~vad
finns det om PEP~8 självt och om stil hos nybörjare?  Motsökningen
redovisas tillsammans med var och en av dem.  \Cref{sec:pep8-metod}
beskriver sökningen, \cref{sec:pep8-resultat} vad den fann,
\cref{sec:pep8-diskussion} vad som begränsar svaret och
\cref{sec:pep8-slutsats} ger det; \cref{sec:pep8-fortsatt} säger vad
som återstår att undersöka, och \cref{tab:sok-pep8} sist i kapitlet
listar de träffar som rör påståendet.

\section{Metod}
\label{sec:pep8-metod}

Frågan söktes ämnesdrivet och tvåsidigt den 6 september 2026 med
verktyget \texttt{scholar} mot de databaser som svarade: OpenAlex (OA),
DBLP, IEEE Xplore (IEEE), arXiv, Web of Science (WoS) och Scopus.
Förkortningen OA står för databasens namn, inte för
\foreignlanguage{english}{open access}.  Semantic Scholar avvisade sin
nyckel och saknas därför helt.  IEEE Xplores dygnskvot tog slut efter de
tjugosex långa frågorna, så de åtta korta frågorna kunde inte ställas
där.  DBLP söker bara i titlar och tar bort citattecken och operatorer,
vilket gör varje flerbegreppsfråga till ett krav på att alla orden står i
samma titel; därför är DBLP-kolumnen noll för de långa frågorna, och
därför ställdes samma frågor en gång till i kort form.

Stödfrågorna söker på de begrepp forskningen använder, inte på
praktikerns ord: \emph{identifier naming}, \emph{program comprehension},
\emph{indentation}, \emph{code layout}, \emph{coding conventions},
\emph{code readability}.  Att ordet \enquote{PEP~8} nästan inte
förekommer i litteraturen är i sig en del av svaret.  Motfrågorna
använder samma begrepp med negativ vokabulär tillagd --- \emph{no
effect}, \emph{no significant difference}, \emph{null result},
\emph{failed replication}, \emph{criticism}, \emph{little evidence}.

\Cref{tab:sok-queries-pep8} visar frågorna och antalet träffar per
databas.  Bokstaven i första kolumnen säger vilken sorts fråga raden är:
S är en stödfråga, C en motfråga och T en kort omfråga ställd för DBLP:s
skull.  Bokstaven i andra kolumnen är delfrågan --- (a)~namn,
(b)~utseende, (c)~konvention, (d)~PEP~8 och nybörjare --- där (e)
betecknar motsökningen.  Frågan står i den booleska form som skickades
till OpenAlex, IEEE Xplore och arXiv; Web of Science fick samma sträng
som \texttt{TS=(\dots)} och Scopus som \texttt{TITLE-ABS-KEY(\dots)}.
Högst 25 träffar hämtades per databas och fråga, så en 25 i tabellen
betyder \enquote{minst 25}, och ett tankstreck i IEEE-kolumnen betyder
att frågan inte prövades där.

De flesta fulltexter fick hämtas utanför databasernas egna länkar, av
skäl som är förlagens: ACM:s och IEEE:s filer är betalspärrade, och i
stället lästes författarnas egna eller lärosätenas kopior.  Sex källor
gick inte att få tag på alls och citeras därför bara för det deras
sammanfattningar själva anger; det står vid var och en av dem i texten.

% Author's note (verification and reproduction):
%   scholar session : variables-pep8
%   exports         : litteratursokning/pep8-stil.csv (classified),
%                     litteratursokning/pep8-stil-full.tex
%   searched        : 2026-09-06; providers openalex dblp ieee arxiv wos
%                     scopus (s2: key rejected HTTP 403; ieee: daily quota
%                     exhausted after the 26 long queries, so the eight
%                     T-rows show 0 for quota reasons).  Limit 25 per
%                     provider and query, so a 25 means "at least 25".
%   classification  : scholar llm context + scholar llm classify
%                     --no-examples, model github_copilot/gpt-5.4; the
%                     low-confidence rows and every qualifying row were
%                     then read by hand.
%   hit-list table  : scholar sessions export variables-pep8 -f table
%                     --bearing-only --lang sv --label sok-pep8
%                     --track "PEP 8 och kodstil"
%                     --theme namngivning=identifierarnamn
%                     --theme layout=layout-och-formatering
%                     --theme konvention=konventionsefterlevnad
%                     --theme pep8=pep8-och-undervisning
%                     -o litteratursokning/pep8-stil
%                     Three edits to the generated fragment: the provider
%                     column heading reads "Databas" rather than the
%                     exporter's "Leverantor" (as in the other decks);
%                     "open access" was wrapped in
%                     \foreignlanguage{english}{}; and the exporter's
%                     paragraph-length caption was cut to one sentence,
%                     its counts and column key moved into the prose of
%                     \cref{sec:pep8-traffar} (review round 4).  Redo all
%                     three after any re-export.
%   provenance      : ltnotes.bib, keys Schankin2018Names
%                     SharifMaletic2010CamelCase Beniamini2017SingleLetter
%                     Scanniello2017Abbreviated Avidan2017VariableNames
%                     Miara1983Indentation Bauer2019Indentation
%                     Hanenberg2024Indentation Oliveira2023FormattingSLR
%                     SantosGerosa2018Readability
%                     Lee2015ConventionViolations
%                     BoogerdMoonen2008CodingStandards
%                     Kirk2024StyleNotQuality
%                     Oliveira2024PythonStyleGuides Kirk2025DistillingPEP8
%                     vanderWerf2024TeachersNaming
%                     Keuning2023CodeQualityMapping
%                     Keuning2017StudentCodeQuality
%   abstract only   : Scanniello2017Abbreviated Hanenberg2024Indentation
%                     BoogerdMoonen2008CodingStandards
%                     Kirk2024StyleNotQuality Kirk2025DistillingPEP8
%                     vanderWerf2024TeachersNaming (six entries; each
%                     says so in its VERIFIED field and in the prose)
%   not obtained    : Binkley et al. 2013 on identifier style (doi
%                     10.1007/s10664-012-9201-4) -- reported through the
%                     systematic review instead; Hofmeister et al.
%                     2017/2019; Ribeiro and Travassos 2018/2023 (the
%                     most worthwhile to chase: it reports four-space
%                     indentation dominating reader preference, which
%                     would qualify dos Santos and Gerosa); Lawrie et al.
%                     2006/2007; Sharafi et al. 2012.  The full dropped
%                     list is in the search report.

\begingroup\footnotesize
\begin{longtable}{@{}ll%
  >{\raggedright\arraybackslash}p{0.40\linewidth}rrrrrr@{}}
  \caption{Sökfrågor och antal träffar per databas för sökningen om
    PEP~8.}%
  \label{tab:sok-queries-pep8}\\
  \toprule
  & & Nyckelord & OA & DBLP & IEEE & arXiv & WoS & Scopus \\
  \midrule
  \endfirsthead
  \multicolumn{9}{@{}l}{\emph{\tablename~\thetable{} (forts.)}}\\
  \toprule
  & & Nyckelord & OA & DBLP & IEEE & arXiv & WoS & Scopus \\
  \midrule
  \endhead
  \midrule \multicolumn{9}{r@{}}{\emph{forts.\ på nästa sida}}\\
  \endfoot
  \bottomrule
  \endlastfoot
  S1 & (a) & \enquote{identifier naming} \textsc{and} (\enquote{program
           comprehension} \textsc{or} readability)
     & 25 & 0 & 15 & 16 & 20 & 25 \\
  S2 & (a) & (\enquote{camelCase} \textsc{or} \enquote{under\_score}
           \textsc{or} \enquote{identifier style}) \textsc{and}
           (comprehension \textsc{or} readability \textsc{or} \enquote{eye
           tracking})
     & 25 & 0 & 5 & 1 & 7 & 9 \\
  S3 & (a) & (\enquote{identifier names} \textsc{or} \enquote{variable
           names}) \textsc{and} \enquote{controlled experiment} \textsc{and}
           comprehension
     & 23 & 0 & 2 & 0 & 5 & 4 \\
  S4 & (a) & (abbreviations \textsc{or} abbreviated \textsc{or}
           \enquote{full words}) \textsc{and} identifiers \textsc{and}
           (comprehension \textsc{or} readability \textsc{or} memory)
     & 25 & 0 & 21 & 25 & 25 & 25 \\
  S5 & (a) & (\enquote{naming conventions} \textsc{or} \enquote{identifier
           quality}) \textsc{and} (defects \textsc{or} faults \textsc{or}
           \enquote{maintenance effort} \textsc{or} maintainability)
     & 25 & 0 & 12 & 14 & 19 & 25 \\
  S6 & (b) & indentation \textsc{and} (\enquote{program comprehension}
           \textsc{or} \enquote{source code} \textsc{or} \enquote{computer
           program}) \textsc{and} (comprehension \textsc{or} readability)
     & 25 & 0 & 5 & 4 & 15 & 17 \\
  S7 & (b) & (\enquote{source code formatting} \textsc{or} \enquote{program
           layout} \textsc{or} \enquote{pretty printing} \textsc{or}
           \enquote{code layout}) \textsc{and} (comprehension \textsc{or}
           readability \textsc{or} maintenance)
     & 25 & 0 & 8 & 4 & 16 & 25 \\
  S8 & (b) & \enquote{code readability} \textsc{and} (model \textsc{or}
           metric \textsc{or} \enquote{human judgement} \textsc{or}
           \enquote{human judgment}) \textsc{and} \enquote{source code}
     & 25 & 0 & 25 & 6 & 25 & 25 \\
  S9 & (b) & (whitespace \textsc{or} \enquote{blank lines} \textsc{or}
           \enquote{line length} \textsc{or} \enquote{visual structure})
           \textsc{and} \enquote{source code} \textsc{and} (comprehension
           \textsc{or} readability)
     & 25 & 0 & 8 & 1 & 5 & 11 \\
  S10 & (c) & (\enquote{coding conventions} \textsc{or} \enquote{coding
           standards} \textsc{or} \enquote{code conventions}) \textsc{and}
           (readability \textsc{or} comprehension \textsc{or}
           maintainability)
     & 25 & 0 & 25 & 25 & 25 & 25 \\
  S11 & (c) & (\enquote{coding style} \textsc{or} \enquote{programming
           style} \textsc{or} \enquote{code style}) \textsc{and}
           (consistency \textsc{or} convention \textsc{or} violation)
           \textsc{and} (comprehension \textsc{or} readability \textsc{or}
           quality)
     & 25 & 0 & 23 & 18 & 25 & 25 \\
  S12 & (c) & (\enquote{code smell} \textsc{or} \enquote{convention
           violation} \textsc{or} \enquote{style violation}) \textsc{and}
           (\enquote{program comprehension} \textsc{or} readability
           \textsc{or} \enquote{developer perception})
     & 25 & 0 & 25 & 21 & 25 & 25 \\
  S13 & (d) & \enquote{PEP 8} \textsc{or} \enquote{PEP8} \textsc{or}
           (\enquote{Python} \textsc{and} \enquote{style guide} \textsc{and}
           \enquote{source code})
     & 25 & 0 & 7 & 7 & 25 & 25 \\
  S14 & (d) & (linter \textsc{or} \enquote{static analysis} \textsc{or}
           \enquote{style checker} \textsc{or} pycodestyle \textsc{or}
           checkstyle) \textsc{and} (\enquote{introductory programming}
           \textsc{or} CS1 \textsc{or} novice \textsc{or} students)
           \textsc{and} (style \textsc{or} feedback)
     & 25 & 0 & 25 & 11 & 25 & 25 \\
  S15 & (d) & (automated \textsc{or} automatic) \textsc{and} (assessment
           \textsc{or} grading \textsc{or} feedback) \textsc{and}
           (\enquote{programming style} \textsc{or} \enquote{code style}
           \textsc{or} \enquote{code quality}) \textsc{and} (students
           \textsc{or} novice)
     & 25 & 0 & 25 & 15 & 25 & 25 \\
  S16 & (d) & (novice \textsc{or} \enquote{introductory programming}
           \textsc{or} CS1) \textsc{and} \enquote{code quality} \textsc{and}
           (style \textsc{or} readability \textsc{or} conventions)
     & 25 & 0 & 23 & 6 & 23 & 25 \\
  T1 & (a) & identifier names comprehension
     & 25 & 2 & -- & 25 & 25 & 25 \\
  T2 & (a) & identifier style eye tracking
     & 25 & 1 & -- & 25 & 3 & 4 \\
  T3 & (b) & indentation comprehension
     & 25 & 4 & -- & 25 & 25 & 25 \\
  T4 & (b) & code readability metric
     & 25 & 5 & -- & 25 & 25 & 25 \\
  T5 & (c) & coding conventions readability
     & 25 & 0 & -- & 25 & 25 & 25 \\
  T6 & (d) & programming style automated feedback students
     & 25 & 0 & -- & 25 & 25 & 24 \\
  \midrule
  C1 & (e) & (\enquote{identifier naming} \textsc{or} \enquote{identifier
           style} \textsc{or} \enquote{variable names}) \textsc{and}
           (\enquote{no significant difference} \textsc{or} \enquote{no
           effect} \textsc{or} \enquote{null result})
     & 25 & 0 & 1 & 25 & 2 & 4 \\
  C2 & (e) & (\enquote{camelCase} \textsc{or} \enquote{under\_score}
           \textsc{or} \enquote{identifier style}) \textsc{and} (replication
           \textsc{or} \enquote{conflicting results} \textsc{or} \enquote{no
           significant})
     & 24 & 0 & 0 & 3 & 5 & 4 \\
  C3 & (e) & indentation \textsc{and} (\enquote{source code} \textsc{or}
           \enquote{computer program}) \textsc{and} (\enquote{no
           significant} \textsc{or} \enquote{no effect} \textsc{or}
           replication \textsc{or} contradictory)
     & 25 & 0 & 1 & 5 & 3 & 6 \\
  C4 & (e) & (\enquote{coding conventions} \textsc{or} \enquote{coding
           standards} \textsc{or} \enquote{code conventions}) \textsc{and}
           (criticism \textsc{or} limitations \textsc{or} \enquote{lack of
           evidence} \textsc{or} \enquote{no effect})
     & 25 & 0 & 25 & 25 & 25 & 25 \\
  C5 & (e) & (\enquote{code readability} \textsc{or} \enquote{program
           comprehension}) \textsc{and} (\enquote{negative results}
           \textsc{or} \enquote{failed replication} \textsc{or}
           \enquote{does not improve} \textsc{or} \enquote{no difference})
     & 25 & 0 & 7 & 25 & 4 & 11 \\
  C6 & (e) & (\enquote{programming style} \textsc{or} \enquote{code style}
           \textsc{or} \enquote{coding style}) \textsc{and} (novice
           \textsc{or} students) \textsc{and} (\enquote{no effect}
           \textsc{or} \enquote{no improvement} \textsc{or} \enquote{not
           associated})
     & 25 & 0 & 0 & 10 & 0 & 1 \\
  C7 & (e) & (\enquote{source code formatting} \textsc{or} \enquote{program
           layout} \textsc{or} \enquote{code layout}) \textsc{and}
           \enquote{controlled experiment} \textsc{and} (\enquote{no
           difference} \textsc{or} \enquote{not significant} \textsc{or}
           inconclusive)
     & 23 & 0 & 0 & 0 & 0 & 1 \\
  C8 & (e) & (linter \textsc{or} \enquote{static analysis} \textsc{or}
           \enquote{style checker}) \textsc{and} (students \textsc{or}
           novice) \textsc{and} (\enquote{no effect} \textsc{or}
           \enquote{did not improve} \textsc{or} limitations \textsc{or}
           \enquote{false positives})
     & 25 & 0 & 25 & 15 & 18 & 25 \\
  C9 & (e) & \enquote{empirical evidence} \textsc{and} (\enquote{coding
           conventions} \textsc{or} \enquote{coding style} \textsc{or}
           \enquote{software engineering practices}) \textsc{and} (weak
           \textsc{or} lacking \textsc{or} inconclusive \textsc{or}
           \enquote{little evidence})
     & 25 & 0 & 0 & 2 & 1 & 1 \\
  C10 & (e) & \enquote{program comprehension} \textsc{and}
           \enquote{controlled experiment} \textsc{and} (replication
           \textsc{and} (\enquote{failed} \textsc{or} \enquote{contrary}
           \textsc{or} \enquote{could not}))
     & 25 & 0 & 0 & 0 & 0 & 1 \\
  T7 & (e) & coding style experiment negative results
     & 25 & 0 & -- & 25 & 19 & 11 \\
  T8 & (e) & program comprehension replication experiment
     & 24 & 2 & -- & 25 & 23 & 21 \\
\end{longtable}
\endgroup

Sammanlagt gav sökningarna 2\,773 träffar, efter sammanslagning av
dubbletter 1\,824 unika poster.  Alla klassades: först
med en språkmodell mot en beskrivning av frågan och fyra kategorier ---
stöder påståendet, kvalificerar eller motsäger det, angränsande delämne,
annat ämne --- och sedan för hand: varje post i de två första
kategorierna lästes.  Arton källor citeras i \cref{sec:pep8-resultat};
tolv av dem lästes i fulltext och sex fanns bara tillgängliga som
sammanfattning, vilket står vid var och en.

\section{Resultat}
\label{sec:pep8-resultat}

\paragraph{(a) Namnen.}  Här finns det starkaste stödet, och den
tydligaste motsägelsen.  Schankin och medförfattare lät utvecklare leta
efter ett betydelsefel i kod med beskrivande respektive korta namn och
fann felet omkring fjorton procent snabbare i den beskrivande
versionen; effekten försvann när uppgiften i stället var att hitta ett
syntaxfel, alltså när koden inte behövde förstås\autocite[avsnitt~5,
7]{Schankin2018Names}.  Samma studie innehåller dock ett förbehåll som
väger tungt just för en nybörjarkurs: \enquote{\foreignlanguage{english}{%
the style of identifier names had a clear impact on program comprehension
for more experienced developers but not for less experienced
developers}}\autocite[sammanfattningen]{Schankin2018Names}.  Studiens
egen sammanfattning av avsnitt~5 är försiktigare: en \emph{måttlig}
verkan hos de vana och en \emph{liten} hos de ovana --- alltså inte
ingen.  Dos Santos och Gerosa lät läsare jämföra kodstycken två
och två och fann att namn byggda av ordboksord upplevdes som läsbarare
med god marginal ($p = 0{,}0001$)%
\autocite[avsnitt~4.2]{SantosGerosa2018Readability}.  Avidan och Feitelson
bytte namn i \emph{verklig} produktionskod och fick samma riktning ---
\enquote{\foreignlanguage{english}{names have a large impact on the
comprehension of code}} --- men fann också något som den här modulen
redan visar med sitt tredje momsprogram: att
\enquote{\foreignlanguage{english}{misleading names, or names that clash
with their types, are worse than meaningless names like consecutive
letters of the alphabet}}\autocite[avsnitt~IX]{Avidan2017VariableNames}.
Ett namn som ljuger är alltså sämre än inget namn alls.

Motsökningen gav två studier som pekar åt andra hållet.  Beniamini och
medförfattare fann i kontrollerade försök \enquote{\foreignlanguage{%
english}{no significant differences in ability to modify the code}} när
bara variabelnamnen byttes ut, och drar slutsatsen att enbokstavsnamn
ibland är att föredra\autocite[sammanfattningen]{Beniamini2017SingleLetter}.
Scanniello och medförfattare jämförde förkortade och fullständiga namn i
ett ursprungsexperiment och tre replikeringar med hundra deltagare och
fann ingen skillnad i att hitta och rätta fel; den studien har bara
kunnat läsas i sammanfattning\autocite{Scanniello2017Abbreviated}.
Skillnaden mellan de två lägren följer uppgiften: att \emph{förstå} kod
går snabbare med beskrivande namn, att \emph{ändra} kod tycks inte göra
det.

Om just den skrivregel modulen lär ut --- små bokstäver och understreck
--- finns en ögonrörelsestudie: understrecksstilen kändes igen snabbare
och med mindre ansträngning än \foreignlanguage{english}{camelCase}, och
nybörjare tjänade dubbelt så mycket på den som vana
programmerare\autocite[avsnitt~VIII]{SharifMaletic2010CamelCase}.  Men
samma studie fann ingen skillnad i träffsäkerhet, och en tidigare studie,
redovisad i den systematiska översikten, fann det omvända för vissa
uppgifter\autocite[avsnitt~4]{Oliveira2023FormattingSLR}.  De två
studier som prövat frågan är alltså oense, och det är översikten som
redovisar att de är det.

\paragraph{(b) Utseendet.}  Den klassiska studien är Miara och
medförfattares experiment från 1983: indrag hjälpte, och det nyttiga
djupet var litet.  \enquote{\foreignlanguage{english}{we conclude that
some indentation does aid program comprehension \dots{} the optimal level
of indentation is 2--4 spaces}}\autocite[avsnitt~5]{Miara1983Indentation}.
Lägg märke till vad det inte säger: två blanksteg gav det högsta
resultatet, inte fyra.  En modern replikering av samma experiment fann
ingen effekt alls --- varken på förståelse, upplevd svårighet eller
ögonrörelser --- och konstaterar samtidigt att det finns
\enquote{\foreignlanguage{english}{little empirical evidence on optimal
indentation depth}}\autocite[avsnitt~8]{Bauer2019Indentation}.  Ett
randomiserat försök från 2024 vänder tillbaka igen: med indragen kod mot
helt oindragen syns en \enquote{\foreignlanguage{english}{strong \dots{}
and large \dots{} effect of indentation}}, en dryg fördubbling av
lästiden utan indrag --- men effekten syns för gruppen, inte för varje
enskild deltagare, och den studien har bara kunnat läsas i
sammanfattning\autocite[sammanfattningen]{Hanenberg2024Indentation}.  De
tre studierna prövar inte samma sak: Miara och Bauer jämför \emph{djup},
Hanenberg och medförfattare jämför indrag mot inget indrag alls.

Dos Santos och Gerosas jämförelser regel för regel ger samma splittrade
bild: att hålla raderna inom åttio tecken gjorde märkbar skillnad
($p = 0{,}0003$), medan fyra blanksteg i indrag inte gjorde det
($p = 0{,}3222$) --- där föredrog till och med 58~procent av läsarna
det kodstycke som bröt mot regeln --- och tomrader mellan
sammanhörande satser gjorde ingen skillnad
alls\autocite[avsnitt~4.2]{SantosGerosa2018Readability}.  Den enda
systematiska översikten på området går igenom femton studier och
tjugosju jämförelser, av vilka sjutton var signifikanta och tio inte, och
sammanfattar läget så här: \enquote{\foreignlanguage{english}{The number
of identified papers, some of which are outdated, and the many null and
contradictory results emphasize the relative lack of work in this
area}}\autocite[sammanfattningen]{Oliveira2023FormattingSLR}.  Bland de
faktorer översikten inte fann några signifikanta resultat för står
mellanslag runt operatorer --- en av de fyra regler den här modulen
räknar upp.

\paragraph{(c) Konventionen som sådan.}  Lee och medförfattare mätte i
210 öppna Javaprojekt hur brott mot Suns stilguide hängde ihop med
läsbarhet och fann ett samband, starkast för indrag, eftersläpande
kommentarer och
dokumentationskommentarer\autocite[avsnitt~7]{Lee2015ConventionViolations}.
Studien är värd att läsa med förbehållet i behåll: läsbarheten mättes
inte på människor utan förutsades med en modell, så resultatet är ett
samband mellan regelbrott och en modells poäng, inte ett belägg för vad
en läsare faktiskt förstår.

Motsökningen gav två invändningar som väger.  Boogerd och Moonen
undersökte MISRA~C, den kodstandard som har mest kontrollmaskineri bakom
sig, och skriver: \enquote{\foreignlanguage{english}{In spite of the
widespread use of coding standards and tools enforcing their rules, there
is little empirical evidence supporting the intuition that they prevent
the introduction of faults in software}} --- och att efterlevnad till och
med kan \emph{öka} antalet fel, eftersom varje ändring kan införa ett
nytt\autocite[sammanfattningen]{BoogerdMoonen2008CodingStandards}.  Den
studien mäter fel, inte läsbarhet, och har bara kunnat läsas i
sammanfattning.  Kirk och medförfattare gör en principiell invändning som
rör undervisning direkt: att blanda ihop stil med kvalitet
\enquote{\foreignlanguage{english}{risks students believing that meeting
guides will \textquote{improve} their code, when in fact it may
not}}\autocite[sammanfattningen]{Kirk2024StyleNotQuality}.  Det är ett
ställningstagande, inte ett experiment, och citeras som ett sådant; också
den har bara kunnat läsas i sammanfattning.

Här finns också ett hål, och det är värt att stanna vid.  Varken
stödfrågorna eller motfrågorna hittade någon studie med mänskliga
deltagare som skiljer \emph{att} koden följer en konvention från
\emph{vilken} konvention den följer.  Argumentet att en gemensam stil
hjälper därför att läsaren vet vad hon ska vänta sig står i inledningen
till flera av de artiklar som redovisats här, men det står inte prövat
någonstans.

\paragraph{(d) PEP~8 självt, och nybörjarna.}  Sökningen fann en enda
studie som prövar PEP~8 direkt, på Python, med nybörjare: Oliveira och
medförfattare lät trettiotvå nybörjare läsa kod med och utan fyra av
PEP~8:s regler och mätte ögonrörelser.  Tre av reglerna höll:
\enquote{\foreignlanguage{english}{non-adherence to PEP8's proper spacing
increases time, number of fixations and horizontal
regressions}}.  Den fjärde gjorde det inte, tvärtom:
\enquote{\foreignlanguage{english}{deviating from PEP8 recommendations
might actually aid in novice programmers' understanding of the
code}}\autocite[avsnitt~8]{Oliveira2024PythonStyleGuides}.  Samma
författare konstaterar i sin inledning att
\enquote{\foreignlanguage{english}{many guides, including PEP8, lack
empirical studies as a foundation}}.  Lägg märke till att just
mellanslagsregeln, som hjälpte här, är en av dem den systematiska
översikten inte fann något signifikant resultat för; studierna använder
olika språk, uppgifter och mått, och de motsäger varandra.

Om PEP~8 som \emph{undervisningstext} finns en studie till, och den är
kritisk: guiden \enquote{\foreignlanguage{english}{[is] not designed for
teaching as [it contains] content that is not suitable for introductory
programming courses and often present[s] advice without explaining why it
should be adopted}}\autocite[sammanfattningen]{Kirk2025DistillingPEP8}
--- också den studien har bara kunnat läsas i sammanfattning.  Det är
precis den avgränsning den här modulen gör: ni ska veta att guiden finns
och följa den, inte lära er den utantill.  Och att ens namnregeln är
självklar i undervisningen stämmer inte: tio intervjuade lärare i
inledande pythonkurser var eniga om att namn ska vara meningsfulla men
oense om vad det betyder, och undervisade sällan om saken
alls\autocite[sammanfattningen]{vanderWerf2024TeachersNaming}.  Även den
intervjustudien har lästs i sammanfattning och ingenting mer.

Om undervisningen finns det mer, men inte om det vi frågar efter.  En
systematisk kartläggning av 195 arbeten om kodkvalitet i utbildning från
1976 till 2022 visar ett fält som till största delen bygger verktyg och
beskriver studenters kod, inte ett som mäter vad stilundervisning ger för
lärande\autocite[avsnitt~5]{Keuning2023CodeQualityMapping}.  Den största
enskilda studien av nybörjares kod --- 2,6 miljoner ögonblicksbilder ---
säger att \enquote{\foreignlanguage{english}{students hardly fix issues
\dots{} and the use of tooling does not have much effect on the
occurrence of issues}}, och råder lärare att välja ut ett litet antal
relevanta kontroller i stället för att slå på
allt\autocite[sammanfattningen, avsnitt~6]{Keuning2017StudentCodeQuality}.

\paragraph{Motsökningen.}  Tolv av de trettiofyra frågorna var
motfrågor, ställda med samma begrepp som stödfrågorna plus negativ
vokabulär, på samma databaser.  De gav 728 av de 2\,773 träffarna och
nästan bara felträffar, och det är i sig lärorikt: ett null-resultat
annonseras sällan i titel eller sammanfattning med orden \emph{no
effect}.  De studier som faktiskt motsäger påståendet kom i stället ur
stödfrågorna och hittades genom läsning --- Bauers replikering heter en
fråga, Beniaminis nollresultat ligger inne i en artikel om
enbokstavsnamn, och Scanniellos fyra experiment är döpta efter uppgiften,
inte efter utfallet.  Motsökningen gav också ett tomrum: ingen studie
hittades som visar att stilundervisning eller stilverktyg förbättrar vad
nybörjare lär sig --- och ingen som visar motsatsen heller.

\section{Diskussion och begränsningar}
\label{sec:pep8-diskussion}

Fem saker begränsar svaret.  För det första är underlaget tunt: hela
frågan om utseendets betydelse vilar på ett femtontal studier, varav den
mest citerade är från 1983.  För det andra handlar nästan ingen av dem om
Python --- de använder Java, C och Pascal --- och den enda som gör det är
från 2024 och prövar fyra regler.  För det tredje mäter studierna olika
saker under samma ord: \emph{läsbarhet} är ibland hur snabbt någon hittar
ett fel, ibland vilket av två kodstycken någon tycker bäst om, ibland vad
en modell räknar fram, och i ett fall antalet fel som smyger sig in.  De
svaren är inte utbytbara, och studier som mäter olika saker kan se ut att
motsäga varandra utan att göra det.  För det fjärde, och viktigast för
den här kursen: Schankins effekt av beskrivande namn syntes tydligt hos
vana utvecklare men bara svagt hos ovana, och nästan alla studier har
yrkesprogrammerare eller senare studenter som deltagare.  Att en regel
hjälper en van läsare är inte ett belägg för att den hjälper en
nybörjare.  För det femte har sex av de arton citerade källorna bara
kunnat läsas i sammanfattning; vid var och en av dem står vad
sammanfattningen säger och ingenting mer.

Till det kommer sökningens egna hål.  Semantic Scholar var otillgängligt
hela tiden och IEEE Xplores kvot tog slut halvvägs.  Att motfrågorna gav
så lite
är dessutom ett metodfynd som gäller vidare än den här frågan: i det här
fältet står nollresultaten inne i artiklar med neutrala titlar, så
tvåsidigheten fick komma ur läsningen av stödträffarna, inte ur den
negativa vokabulären.

\section{Slutsats}
\label{sec:pep8-slutsats}

Delvis, och svagare än man skulle kunna tro.  Det som har stöd är att koden
skrivs \emph{för en läsare}: beskrivande namn gör att en läsare hittar
ett fel snabbare, ett namn som ljuger är sämre än inget namn alls, korta
rader gör kod märkbart läsbarare, och indragen kod läses ungefär dubbelt
så snabbt som helt oindragen.  Det som inte har stöd är guiden som
helhet, och flera av just de regler den här modulen räknar upp: fyra
blanksteg i indrag har prövats och gav ingen skillnad --- en knapp
majoritet av rösterna gick rentav till den kod som bröt mot regeln ---
och mellanslag runt
operatorer gav inget signifikant resultat i den systematiska översikten,
fast det hjälpte i den enda studien på pythonnybörjare.  Den studien fann
att tre av fyra prövade PEP~8-regler hjälpte och att den fjärde stjälpte.

Att låta bli helt är inte heller vad underlaget rekommenderar; ingenting
i någondera riktningen säger att stil inte spelar roll.  Men på frågan om
det \emph{kostar} något att låta bli finns inget svar: ingen studie visar
att stilundervisning gör nybörjare till bättre programmerare, en
industristudie av kodstandarder finner tvärtom lite belägg för att de
förhindrar fel, och den enda studien av PEP~8 som undervisningstext säger
att guiden inte är skriven för att undervisas ur.

Modulen håller därför kvar PEP~8, men av det skäl underlaget bär.  De
regler den lutar sig hårdast mot --- beskrivande namn och en genomtänkt
form --- är också de mest prövade.  Att skriva fyra blanksteg i indrag är
däremot ingen vetenskapligt belagd förbättring; det är att göra som alla
andra i Pythongemenskapen, vilket är ett annat och ärligare skäl.  Och
den viktigaste lärdomen för en kurs finns i den största studien av
nybörjares kod: peka en stilkontroll mot ett litet, valt antal regler,
inte mot allt ett verktyg kan klaga
på\autocite[avsnitt~6]{Keuning2017StudentCodeQuality}.

\section{Fortsatt arbete}
\label{sec:pep8-fortsatt}

Sökningen är inte färdig.  \Cref{sec:pep8-diskussion} räknar upp vad som
brister; det här avsnittet säger vad som ska göras åt det.  De tomma
databascellerna ska fyllas en annan dag: samtliga frågor en gång till på
Semantic Scholar, och de åtta korta frågorna på IEEE Xplore, vars
tankstreck är en kvot som tog slut och inte ett utfall.  Närmare
sexhundra av de
unika posterna saknade sammanfattning när klassningen gjordes, så för
dem vilar gallringen på titeln allena; posterna ska kompletteras och
klassas om.  Och de sex källor som bara är lästa i sammanfattning ska
läsas i fulltext\autocite{Scanniello2017Abbreviated,
  Hanenberg2024Indentation,BoogerdMoonen2008CodingStandards,
  Kirk2024StyleNotQuality,Kirk2025DistillingPEP8,
  vanderWerf2024TeachersNaming} --- först det randomiserade
indragsförsöket, som bär kapitlets starkaste resultat om utseendet och
ändå citeras ur sin egen sammanfattning.

Fyra arbeten kom fram i sökningen men gick inte att få tag på, och vart
och ett skulle avgöra något bestämt.  Ribeiro och Travassos artikel i
\emph{CLEI Electronic Journal} 2018, och dess utvidgning i
\emph{Empirical Software Engineering} 2023, handlar just om de
motsägelser kapitlet redovisar, och sidan som beskriver dem uppger att
fyra blanksteg i indrag dominerade läsarnas preferens --- vilket skulle
kvalificera dos Santos och Gerosas resultat åt andra hållet.  Saliba med
fleras studie från SIGCSE 2024 är det närmaste en klassrumsstudie av om
stilåterkoppling förbättrar studenters stil, alltså precis den lucka
delfråga~(d) lämnar öppen.  Binkley med fleras experiment från 2013 är
den studie som motsäger Sharif och Maletic om
\foreignlanguage{english}{camelCase} mot understreck, och redovisas här
i andra hand genom den systematiska översikten.  Och Hofmeister,
Siegmund och Holt prövade korta mot fullständiga identifierarnamn på 72
yrkesutvecklare, vilket skulle pröva namnfyndet i en andra
population.\footnote{Ribeiro och Travassos: \textsc{doi}
  10.19153/cleiej.21.1.5 och 10.1007/s10664-023-10360-5.  Saliba med
  flera: 10.1145/3626252.3630889.  Binkley med flera:
  10.1007/s10664-012-9201-4.  Hofmeister med flera:
  10.1109/SANER.2017.7884623 och 10.1007/s10664-018-9621-x.}

Vad som saknas för ett säkrare svar är framför allt studier på dem
frågan gäller.  Nästan alla experiment i kapitlet har yrkesprogrammerare
eller studenter längre fram i utbildningen som deltagare, och den enda
som prövar PEP~8 på pythonnybörjare prövar fyra regler på trettiotvå
personer.  Det som skulle avgöra saken är ett kontrollerat försök där
nybörjare läser samma program med och utan var och en av just de regler
ett material som det här lär ut, med förståelse som utfall och inte
tycke, och därtill en replikering av jämförelsen regel för regel i en
annan population.  Två delfrågor är dessutom obesvarade i hela
underlaget: om det spelar roll \emph{att} koden följer en konvention
skilt från \emph{vilken} konvention den följer, och om stilen har någon
verkan på \emph{underhållet} --- allt som redovisas här mäter läsning,
inte vad det kostar att ändra koden ett halvår senare.  Utfallet skulle
synas i materialet: visade ett sådant försök att reglerna hjälper också
nybörjare, kunde \cref{sec:kommentarer} åter säga att stilreglerna finns
för läsarens skull, och visade det motsatsen skulle skälet krympa till
det som står där nu --- att göra som alla andra i Pythongemenskapen.

\section{Träffar som rör påståendet}
\label{sec:pep8-traffar}

\Cref{tab:sok-pep8} listar de poster som rör påståendet: arton citerade,
83 som stöder det och 70 som kvalificerar eller motsäger det, av de
1\,824 unika träffarna.  De 631 angränsande och de 1\,022 felträffarna
listas inte.  Skälet till varje beslut står i tabellens sista kolumn, och
för de rader som en språkmodell klassade anges modellens konfidens.
Databasernas förkortningar är desamma som i \cref{tab:sok-queries-pep8},
där också antalet träffar per fråga står.
\begingroup\footnotesize
\begin{longtable}{@{}%
  >{\raggedright\arraybackslash}p{0.44\textwidth}c%
  >{\raggedright\arraybackslash}p{0.13\textwidth}%
  >{\raggedright\arraybackslash}p{0.26\textwidth}@{}}
\caption{Träffar som rör påståendet om PEP 8 och kodstil.}\label{tab:sok-pep8}\\
\toprule
Titel & År & Databas & Skäl \\
\midrule
\endfirsthead
\multicolumn{4}{@{}l}{\emph{\tablename~\thetable{} (forts.)}}\\
\toprule
Titel & År & Databas & Skäl \\
\midrule
\endhead
\midrule \multicolumn{4}{r@{}}{\emph{forts.\ på nästa sida}}\\
\endfoot
\bottomrule
\endlastfoot
A Systematic Literature Review on the Impact of Formatting Elements on Code Legibility & 2022 & arXiv, OA & \emph{citerad}: layout och formatering \\
A Systematic Mapping Study of Code Quality in Education & 2023 & OA & \emph{citerad}: PEP 8 och undervisning \\
An eye tracking study on camelcase and under-score identifier styles & 2010 & Scopus & \emph{citerad}: namn på variabler och andra identifierare \\
Assessing Python Style Guides: An Eye-Tracking Study with Novice Developers & 2024 & arXiv & \emph{citerad}: PEP 8 och undervisning \\
Assessing the value of coding standards: An empirical study & 2008 & OA & \emph{citerad}: att följa en konvention \\
Code Quality Issues in Student Programs & 2017 & OA & \emph{citerad}: PEP 8 och undervisning \\
Code Style != Code Quality & 2024 & Scopus & \emph{citerad}: PEP 8 och undervisning \\
Descriptive Compound Identifier Names Improve Source Code Comprehension & 2018 & IEEE, DBLP & \emph{citerad}: namn på variabler och andra identifierare \\
Distilling PEP 8 for Teaching Introductory Programming & 2025 & Scopus & \emph{citerad}: PEP 8 och undervisning \\
Effect analysis of coding convention violations on readability of post-delivered code & 2015 & Scopus & \emph{citerad}: att följa en konvention \\
Effects of Variable Names on Comprehension: An Empirical Study & 2017 & Scopus & \emph{citerad}: namn på variabler och andra identifierare \\
Fixing Faults in C and Java Source Code: Abbreviated vs. Full-Word Identifier Names & 2017 & WoS & \emph{citerad}: namn på variabler och andra identifierare \\
Impacts of coding practices on readability & 2018 & Scopus & \emph{citerad}: layout och formatering \\
Indentation and reading time: a randomized control trial on the differences between generated indented and non-indented if-statements & 2024 & Scopus & \emph{citerad}: layout och formatering \\
Indentation: Simply a matter of style or support for program comprehension? & 2019 & Scopus & \emph{citerad}: layout och formatering \\
Meaningful Identifier Names: The Case of Single-Letter Variables & 2017 & OA & \emph{citerad}: namn på variabler och andra identifierare \\
Program indentation and comprehensibility & 1983 & OA & \emph{citerad}: layout och formatering \\
Teachers' Beliefs and Practices on the Naming of Variables in Introductory Python Programming Courses & 2024 & WoS & \emph{citerad}: PEP 8 och undervisning \\
A metric for software readability & 2008 & OA & stöder påståendet (0.99) \\
A Resource to Support Novices Refactoring Conditional Statements & 2022 & Scopus & stöder påståendet (0.78) \\
A Static Analysis Tool in CS1: Student Usage and Perceptions of PythonTA & 2024 & OA & stöder påståendet (0.78) \\
A study of different coding styles affecting code readability & 2013 & Scopus & stöder påståendet (0.82) \\
A Tutoring System to Learn Code Refactoring & 2021 & WoS & stöder påståendet (0.84) \\
Acknowledging Good Java Code with Code Perfumes & 2024 & arXiv, WoS & stöder påståendet (0.87) \\
AI Teaches the Art of Elegant Coding: Timely, Fair, and Helpful Style Feedback in a Global Course & 2024 & OA & stöder påståendet (0.92) \\
An Empirical Investigation on Source Code Comprehension Performance Between Naming and Structural Issues with and without Generative AI Assistance & 2026 & Scopus & stöder påståendet (0.71) \\
An Empirical Study of the Relationships between Code Readability and Software Complexity & 2019 & arXiv & stöder påståendet (0.88) \\
An eye tracking study assessing source code readability rules for program comprehension & 2024 & WoS & stöder påståendet (0.99) \\
An Interactive Tool to Improve Program Readability for Novice Students & 2024 & Scopus & stöder påståendet (0.74) \\
Analysis of Learning Behavior in an Automated Programming Assessment Environment: A Code Quality Perspective & 2020 & OA & stöder påståendet (0.80) \\
Asking Novices to Evaluate Code Quality: Understandability and Efficiency & 2025 & Scopus & stöder påståendet (0.78) \\
Assisting Novice Developers Learning in Flutter Through Cognitive-Driven Development & 2024 & arXiv & stöder påståendet (0.81) \\
Automated Feedback for Quality Assurance in Software Engineering Education & 2010 & IEEE & stöder påståendet (0.74) \\
Automated style feedback for advanced beginner Java programmers & 2016 & IEEE & stöder påståendet (0.83) \\
Automatic Segmentation of Method Code into Meaningful Blocks to Improve Readability & 2011 & IEEE & stöder påståendet (0.97) \\
Brain activity measurement during program comprehension with NIRS & 2014 & OA & stöder påståendet (0.90) \\
Can students help themselves? An investigation of students' feedback on the quality of the source code & 2018 & IEEE & stöder påståendet (0.80) \\
Code Readability Testing, an Empirical Study & 2016 & OA & stöder påståendet (0.90) \\
Developers talking about code quality & 2023 & OA & stöder påståendet (0.88) \\
Developing students' code style skills & 2023 & Scopus & stöder påståendet (0.80) \\
Does ChatGPT Help Novice Programmers Write Better Code? Results From Static Code Analysis & 2024 & IEEE & stöder påståendet (0.90) \\
Does the \textquotedbl{}Refactor to Understand\textquotedbl{} Reverse Engineering Pattern Improve Program Comprehension? & 2005 & OA & stöder påståendet (0.95) \\
Does the discipline of preprocessor annotations matter? A controlled experiment & 2013 & Scopus & stöder påståendet (0.76) \\
Does the modern code inspection have value? & 2001 & IEEE & stöder påståendet (0.92) \\
Eastwood-Tidy: C Linting for Automated Code Style Assessment in Programming Courses & 2023 & WoS & stöder påståendet (0.86) \\
Effective identifier names for comprehension and memory. & 2007 & DBLP & stöder påståendet (0.81) \\
Effects of Human vs. Automatic Feedback on Students' Understanding of AI Concepts and Programming Style & 2020 & arXiv & stöder påståendet (0.77) \\
Enhancing Practice and Achievement in Introductory Programming With a Robot Olympics & 2015 & IEEE & stöder påståendet (0.76) \\
Experiences with TA-Bot in CS1 & 2023 & OA & stöder påståendet (0.82) \\
Exploring Code Comprehension in Scientific Programming: Preliminary Insights from Research Scientists & 2025 & arXiv & stöder påståendet (0.92) \\
Exploring CS1 Student's Notions of Code Quality & 2023 & Scopus & stöder påståendet (0.82) \\
Exploring the Influence of Identifier Names on Code Quality: An Empirical Study & 2010 & OA & stöder påståendet (0.92) \\
Eye-Tracking Insights into the Effects of Type Annotations and Identifier Naming & 2026 & Scopus & stöder påståendet (0.86) \\
Fixing Your Own Smells: Adding a Mistake-Based Familiarisation Step When Teaching Code Refactoring & 2024 & Scopus & stöder påståendet (0.78) \\
Foobaz: Variable Name Feedback for Student Code at Scale & 2015 & WoS & stöder påståendet (0.90) \\
Forming Groups for Collaborative Learning in Introductory Computer Programming Courses Based on Students' Programming Styles: An Empirical Study & 2006 & IEEE & stöder påståendet (0.83) \\
From Obfuscation to Comprehension & 2015 & Scopus & stöder påståendet \\
How Automated Feedback is Delivered Matters: Formative Feedback and Knowledge & 2019 & WoS & stöder påståendet (0.77) \\
How do Developers Improve Code Readability? An Empirical Study of Pull Requests & 2023 & OA, arXiv & stöder påståendet (0.86) \\
How students choose names: A replication study & 2022 & Scopus & stöder påståendet (0.81) \\
Impact of maintainability defects on code inspections & 2010 & Scopus & stöder påståendet (0.70) \\
Improving Code Readability Models with Textual Features & 2016 & WoS & stöder påståendet (0.98) \\
Improving code readability models with textual features & 2016 & Scopus & stöder påståendet (0.98) \\
Improving computer program readability to aid modification & 1982 & OA & stöder påståendet (0.94) \\
Improving the Readability of Automatically Generated Tests using Large Language Models & 2024 & arXiv & stöder påståendet (0.93) \\
Indentation and Reading Time: A Controlled Experiment on the Differences Between Generated Indented and Non-indented JSON Objects & 2024 & Scopus & stöder påståendet (0.84) \\
Indentation, documentation and programmer comprehension. & 1982 & DBLP & stöder påståendet (0.83) \\
Indentation: A Simple Matter of Style or Support for Program Comprehension? & 2018 & DBLP & stöder påståendet (0.84) \\
Introducing Code Quality at CS1 Level: Examples and Activities & 2025 & Scopus & stöder påståendet (0.80) \\
Introducing Code Quality in the CS1 Classroom & 2024 & Scopus & stöder påståendet (0.79) \\
Investigating Naming Convention Adherence in Java References & 2015 & WoS & stöder påståendet (0.84) \\
Is \textquotedbl{}notDone\textquotedbl{} the Same as \textquotedbl{}!done\textquotedbl{}? The Effect of Different Ways for Expressing Negation & 2025 & WoS & stöder påståendet (0.94) \\
It's Never too Early to Learn About Code Quality & 2023 & OA & stöder påståendet (0.84) \\
Learning with Style: Improving Student Code-Style Through Better Automated Feedback & 2024 & Scopus & stöder påståendet (0.78) \\
Lesion analysis of the brain areas involved in language comprehension & 2004 & OA & stöder påståendet (0.97) \\
Maintainability and Source Code Conventions: An Analysis of Open Source Projects & 2011 & OA & stöder påståendet (0.81) \\
Naming Practices in Java Projects: An Empirical Study & 2021 & WoS & stöder påståendet (0.82) \\
On method ordering & 2016 & IEEE, WoS & stöder påståendet (0.95) \\
On the Use of Static Analysis to Engage Students with Software Quality Improvement: An Experience with PMD & 2023 & arXiv & stöder påståendet (0.78) \\
Promoting Code Quality via Automated Feedback on Student Submissions & 2021 & IEEE & stöder påståendet (0.83) \\
PROPER PROGRAMMING STYLE - RIGHT FROM THE BEGINNING & 2019 & WoS & stöder påståendet (0.73) \\
Reordering Program Statements for Improving Readability & 2013 & IEEE & stöder påståendet (0.95) \\
Shorter Identifier Names Take Longer to Comprehend & 2017 & WoS, OA & stöder påståendet (0.99) \\
SOBO: A Feedback Bot to Nudge Code Quality in Programming Courses & 2023 & arXiv & stöder påståendet (0.77) \\
SOBO: A Feedback Bot to Nudge Code Quality in Programming Courses & 2025 & Scopus & stöder påståendet (0.76) \\
Teaching Students to Recognize and Implement Good Coding Style & 2017 & OA & stöder påståendet (0.90) \\
THE EFFECT OF INDENTATION ON PROGRAM COMPREHENSION & 1984 & WoS, DBLP & stöder påståendet (0.91) \\
The impact of identifier style on effort and comprehension & 2013 & Scopus & stöder påståendet (0.90) \\
The Role of Source Code Vocabulary in Programming Teaching and Learning & 2020 & WoS, IEEE & stöder påståendet (0.93) \\
To camelcase or under-score & 2009 & Scopus & stöder påståendet (0.80) \\
To CamelCase or under\_score & 2019 & IEEE, WoS & stöder påståendet (0.96) \\
Tool assisted identifier naming for improved software readability: an empirical study & 2005 & IEEE, WoS & stöder påståendet (0.99) \\
Tracing vs. comprehension: On different levels of understanding Boolean expressions & 2026 & OA & stöder påståendet (0.78) \\
Typographic Style is More than Cosmetic & 1990 & Scopus & stöder påståendet (0.87) \\
Uncovering the Novice Experience: A Think-Aloud Study of Student Interactions with an Automated Coding Style Tool & 2026 & Scopus & stöder påståendet (0.79) \\
Understanding Test Convention Consistency as a Dimension of Test Quality & 2024 & WoS & stöder påståendet (0.91) \\
Unlocking Student Potential With TA-Bot: Timely Submissions and Improved Code Style & 2025 & Scopus & stöder påståendet (0.78) \\
Utilizing programming traces to explore and model the dimensions of novices' codewriting skill & 2023 & OA & stöder påståendet (0.76) \\
What's in a Display Name? An Empirical Study on the Use of Display Names in Open-Source JUnit Tests & 2024 & WoS & stöder påståendet (0.76) \\
Whats in a Name? A Study of Identifiers & 2006 & OA & stöder påståendet (0.99) \\
Work-In-Progress: Code Quality Issues of Computing Undergraduates & 2022 & IEEE & stöder påståendet (0.73) \\
40 Years of Designing Code Comprehension Experiments: A Systematic Mapping Study & 2023 & OA & kvalificerar eller motsäger påståendet (0.78) \\
\textquotedbl{}I know it when I see it\textquotedbl{} Perceptions of Code Quality & 2018 & OA & kvalificerar eller motsäger påståendet (0.74) \\
A Large Scale Empirical Study of the Impact of Spaghetti Code and Blob Anti-patterns on Program Comprehension & 2020 & arXiv & kvalificerar eller motsäger påståendet (0.68) \\
A large-scale comparative analysis of Coding Standard conformance in Open-Source Data Science projects & 2020 & Scopus & kvalificerar eller motsäger påståendet (0.72) \\
A Literature-Informed Model for Code Style Principles to Support Teachers of Text-Based Programming & 2024 & OA & kvalificerar eller motsäger påståendet (0.90) \\
A paradigm for programming style research & 1988 & OA & kvalificerar eller motsäger påståendet (0.95) \\
A Retrospective Study to Understand Student Coding Style Quality via Static Analysis & 2025 & Scopus & kvalificerar eller motsäger påståendet (0.88) \\
A simpler model of software readability & 2011 & OA & kvalificerar eller motsäger påståendet (0.96) \\
An eye tracking study assessing the impact of background styling in code editors on novice programmers' code understanding & 2023 & Scopus & kvalificerar eller motsäger påståendet (0.84) \\
Assessing Consensus of Developers' Views on Code Readability & 2024 & arXiv & kvalificerar eller motsäger påståendet (0.86) \\
Attributes Influencing the Reading and Comprehension of Source Code -- Discussing Contradictory Evidence & 2018 & OA & kvalificerar eller motsäger påståendet (0.98) \\
Automatic Programming Assessment System for a Computer Science Bridge Course - An Experience Report & 2022 & IEEE & kvalificerar eller motsäger påståendet (0.78) \\
Catalog of Code Quality Defects in Introductory Programming & 2024 & Scopus & kvalificerar eller motsäger påståendet (0.85) \\
Code Convention Adherence in Research Data Infrastructure Software: An Exploratory Study & 2019 & IEEE & kvalificerar eller motsäger påståendet (0.78) \\
Code Quality Defects in Introductory Programming & 2025 & OA & kvalificerar eller motsäger påståendet (0.93) \\
Code Readability in the Age of Large Language Models: An Industrial Case Study from Atlassian & 2025 & arXiv & kvalificerar eller motsäger påståendet (0.84) \\
Common Code Quality Issues of Novice Java Programmers: A Comprehensive Analysis of Student Assignments & 2023 & Scopus & kvalificerar eller motsäger påståendet (0.82) \\
Dealing with Faults in Source Code: Abbreviated vs. Full-Word Identifier Names & 2013 & OA & kvalificerar eller motsäger påståendet (0.99) \\
Does code review really remove coding convention violations? & 2020 & OA & kvalificerar eller motsäger påståendet (0.93) \\
Does Code Structure Affect Comprehension? On Using and Naming Intermediate Variables & 2021 & arXiv & kvalificerar eller motsäger påståendet (0.97) \\
Effect of Range Naming Conventions on Reliability and Development Time for Simple Spreadsheet Formulas & 2011 & arXiv & kvalificerar eller motsäger påståendet (0.95) \\
Effects of Cloned Code on Software Maintainability: A Replicated Developer Study & 2013 & WoS & kvalificerar eller motsäger påståendet (0.93) \\
Effects of cloned code on software maintainability: A replicated developer study & 2013 & Scopus & kvalificerar eller motsäger påståendet (0.92) \\
Evaluating Code Readability and Legibility: An Examination of Human-centric Studies & 2020 & Scopus & kvalificerar eller motsäger påståendet (0.81) \\
Evaluating Programming Performance with Keystroke Characteristics: An Empirical Experiment & 2013 & IEEE & kvalificerar eller motsäger påståendet (0.84) \\
Evaluating the relation between coding standard violations and faultswithin and across software versions & 2009 & OA & kvalificerar eller motsäger påståendet (0.96) \\
Experimental investigations of the utility of detailed flowcharts in programming & 1977 & OA & kvalificerar eller motsäger påståendet (0.94) \\
Exploring Software Measures to Assess Program Comprehension & 2011 & OA & kvalificerar eller motsäger påståendet (0.73) \\
Eye Tracking Analysis of Code Layout, Crowding and Dyslexia - An Open Data Set & 2021 & Scopus & kvalificerar eller motsäger påståendet (0.78) \\
Eye tracking analysis of computer program comprehension in programmers with dyslexia & 2018 & OA & kvalificerar eller motsäger påståendet (0.86) \\
Framework for measuring program comprehension & 2012 & OA & kvalificerar eller motsäger påståendet (0.77) \\
From Code Complexity Metrics to Program Comprehension & 2023 & OA & kvalificerar eller motsäger påståendet (0.70) \\
From Code Smells to Curriculum Interventions: Clean Code Deficiencies in Undergraduate Software Engineering Projects & 2026 & Scopus & kvalificerar eller motsäger påståendet (0.74) \\
How do annotations affect Java code readability? & 2024 & arXiv & kvalificerar eller motsäger påståendet (0.96) \\
How Readable is Model-generated Code? Examining Readability and Visual Inspection of GitHub Copilot & 2022 & arXiv & kvalificerar eller motsäger påståendet (0.89) \\
Human vs. Automated Coding Style Grading in Computing Education & 2020 & OA & kvalificerar eller motsäger påståendet (0.94) \\
Identifier Names in Computer Programs: Literature Review & 2023 & Scopus & kvalificerar eller motsäger påståendet (0.86) \\
Identifier Names, Comprehension and Code Metrics: Literature Review & 2024 & Scopus & kvalificerar eller motsäger påståendet (0.87) \\
Impact of AI Debugging on Student Code Complexity and Readability: A Literature Review & 2025 & IEEE & kvalificerar eller motsäger påståendet (0.83) \\
Impact of Limited Memory Resources & 2008 & IEEE & kvalificerar eller motsäger påståendet (0.99) \\
Impacts of Coding Practices on Readability & 2018 & IEEE & kvalificerar eller motsäger påståendet (0.98) \\
Improving Code: The (Mis) Perception of Quality Metrics & 2018 & OA & kvalificerar eller motsäger påståendet (0.95) \\
Improving Source Code Readability: Theory and Practice & 2019 & OA & kvalificerar eller motsäger påståendet (0.86) \\
Inference-based and expectation-based processing in program comprehension & 2001 & IEEE & kvalificerar eller motsäger påståendet (0.71) \\
Investigating the Impact of Automated Code Quality Feedback in an Embedded Systems Course & 2025 & Scopus & kvalificerar eller motsäger påståendet (0.77) \\
Investigating the Readability of Test Code: Combining Scientific and Practical Views & 2024 & arXiv & kvalificerar eller motsäger påståendet (0.83) \\
Mnemonic Variable Names in Parsons Puzzles & 2019 & OA & kvalificerar eller motsäger påståendet (0.98) \\
On Assuring Learning About Code Quality & 2020 & WoS & kvalificerar eller motsäger påståendet (0.90) \\
On the Importance and Shortcomings of Code Readability Metrics: A Case Study on Reactive Programming. & 2021 & DBLP, arXiv & kvalificerar eller motsäger påståendet (0.72) \\
On the Investigation of Empirical Contradictions - Aggregated Results of Local Studies on Readability and Comprehensibility of Source Code & 2023 & Scopus & kvalificerar eller motsäger påståendet (0.86) \\
Personalization of Code Readability Evaluation Based on LLM Using Collaborative Filtering & 2024 & arXiv & kvalificerar eller motsäger påståendet (0.78) \\
Personalized Code Readability Assessment: Are We There Yet? & 2025 & arXiv & kvalificerar eller motsäger påståendet (0.88) \\
Program Comprehension Does Not Primarily Rely On the Language Centers of the Human Brain & 2023 & OA & kvalificerar eller motsäger påståendet (0.78) \\
Pythonic vs Refactorable Pythonic: On the relationship between Pythonic idioms and code quality in machine learning projects & 2026 & WoS & kvalificerar eller motsäger påståendet (0.83) \\
Static Analyses in Python Programming Courses & 2019 & OA & kvalificerar eller motsäger påståendet (0.80) \\
Teacher Perceptions of Feedback in High School Programming Education: A Thematic Analysis & 2020 & WoS & kvalificerar eller motsäger påståendet (0.82) \\
The Effect of Crowding on the Reading of Program Code for Programmers with Dyslexia & 2021 & Scopus & kvalificerar eller motsäger påståendet (0.88) \\
The Effect of Information Content and Length on Name Recollection & 2022 & WoS & kvalificerar eller motsäger påståendet (0.97) \\
The Prevalence of Code Smells in Machine Learning projects & 2021 & IEEE, arXiv & kvalificerar eller motsäger påståendet (0.82) \\
The role of experience and ability in comprehension tasks supported by UML stereotypes & 2007 & WoS & kvalificerar eller motsäger påståendet (0.72) \\
The role of experience and ability in comprehension tasks supported by UML stereotypes & 2007 & Scopus & kvalificerar eller motsäger påståendet (0.71) \\
To What Extent Cognitive-Driven Development Improves Code Readability? & 2022 & arXiv & kvalificerar eller motsäger påståendet (0.93) \\
Towards a Cognitive Model of Dynamic Debugging: Does Identifier Construction Matter? & 2024 & WoS & kvalificerar eller motsäger påståendet (0.76) \\
Towards Comprehensive Assessment of Code Quality at CS1-Level: Tools, Rubrics and Refactoring Rules & 2024 & IEEE & kvalificerar eller motsäger påståendet (0.87) \\
Towards Inclusive Source Code Readability Based on the Preferences of Programmers with Visual Impairments & 2024 & WoS & kvalificerar eller motsäger påståendet (0.96) \\
Who Introduces and Who Fixes? Analyzing Code Quality in Collaborative Student's Projects & 2025 & arXiv & kvalificerar eller motsäger påståendet (0.86) \\
Who Introduces and Who Fixes? Code Quality Dynamics in Collaborative Embedded Systems Projects & 2025 & WoS & kvalificerar eller motsäger påståendet (0.84) \\
Who is right? Evaluating empirical contradictions in the readability \& comprehensibility of source code & 2017 & Scopus & kvalificerar eller motsäger påståendet (0.90) \\
WIP: Exploratory Analysis of Code Quality Issues in CS1 & 2025 & Scopus & kvalificerar eller motsäger påståendet (0.74) \\
Women and men - Different but equal: On the impact of identifier style on source code reading & 2012 & Scopus & kvalificerar eller motsäger påståendet (0.90) \\
\end{longtable}
\endgroup

